\paragraph{后文使用的实分析事实。}
若 $v$ 绝对连续，则对每个常数 $\alpha$，有
\begin{equation}
 v'=0\quad\text{几乎处处于 }\{v=\alpha\}.
 \label{eq:ac-level-set}
\end{equation}
特别地，两个绝对连续函数在可测集合上相等，则其导数在该集合上几乎处处相等。
这是绝对连续函数几乎处处可微与水平集的聚点性质的直接推论\citep[\S222K、\S225C(b)]{Fremlin2001Volume2}。
若 $F$ 在集合 $E\subset\R^d$ 上相对局部 Lipschitz，$x:[0,T]\to E$ 绝对连续且 $T<\infty$，
则 $F\circ x$ 绝对连续；可参照 Lipschitz 延拓与复合的标准结论\citep[\S262B(b)、\S262X(l)]{Fremlin2001Volume2}。
这里的相对局部、向量轨道版本由像集紧性作有限局部化，
在每段使用 $\sum_k|F(x(b_k))-F(x(a_k))|\le L\sum_k\|x(b_k)-x(a_k)\|$ 即得。
在 $F$ 具有环境开邻域上的 $C^1$ 局部表达式处，沿轨道使用通常链式法则；
容量边界和安全时间集仍需下文分别处理，不能用状态空间中的几乎处处可微代替沿轨道论证。
